\documentclass[11pt]{article}

\usepackage[T1]{fontenc}
\usepackage{lmodern}
\usepackage{amsmath,amssymb,mathtools}
\usepackage{geometry}
\usepackage{microtype}
\usepackage{hyperref}
\usepackage{enumitem}

\geometry{margin=1in}
\hypersetup{
  colorlinks=true,
  linkcolor=blue,
  urlcolor=blue,
  citecolor=blue
}

\newcommand{\CP}{\mathbb{CP}}
\newcommand{\RP}{\mathbb{RP}}
\newcommand{\RR}{\mathbb{R}}
\newcommand{\CC}{\mathbb{C}}
\newcommand{\HH}{\mathbb{H}}
\newcommand{\mzero}{\mathfrak m_0}
\newcommand{\mthree}{\mathfrak m_3}
\newcommand{\su}{\mathfrak{su}}
\newcommand{\so}{\mathfrak{so}}
\newcommand{\Tr}{\operatorname{Tr}}

\title{The \texorpdfstring{$\mathfrak m_3$}{m3} Doublet and Custodial
\texorpdfstring{$\rho=1$}{rho=1} in the
\texorpdfstring{$\CP^2$}{CP2}/Aloff--Wallach Kaluza--Klein Carrier}
\author{KK Orchestra working lecture}
\date{June 1, 2026}

\begin{document}

\maketitle

\begin{abstract}
This lecture isolates the clean custodial part of the
\(\CP^2\)/Aloff--Wallach programme.  The internal seven-manifold
\[
SO(3)\longrightarrow X_{1,1}=SU(3)/U(1)_{1,1}\longrightarrow \CP^2
\]
has a tangent split
\[
\mathfrak m=\mzero\oplus\mthree,\qquad
\dim_{\RR}\mzero=3,\qquad \dim_{\RR}\mthree=4 .
\]
The four-dimensional block \(\mthree\simeq \CC^2\simeq\RR^4\) supplies the
electroweak Higgs doublet.  Since \(\mthree\) is a single irreducible real
\(U(2)\)-module, an invariant quadratic form restricts to one scalar on
\(\mthree\).  A vacuum in this single doublet block gives
\(\kappa_c=\kappa_n\) and hence the tree-level custodial relation
\(\rho=1\).  The vertical \(\mzero\) triplet, D6 branch terms, and De Vries
spectral questions are separate consistency layers.
\end{abstract}

\tableofcontents

\section{Reading Map}

The lecture has two audiences in mind.

\begin{itemize}[leftmargin=2em]
  \item For a physicist, the central calculation is the electroweak
  mass matrix of one Higgs doublet, written with the internal metric factor
  carried along.
  \item For a geometer, the central point is the isotropy representation
  of \(X_{1,1}\).  The block \(\mthree\) is a real irreducible four-plane,
  while \(\mzero\) is the vertical three-plane.
\end{itemize}

The conclusion is deliberately modest.  The \(\rho=1\) result is a
standard custodial statement with a geometric origin.  It does not derive a
De Vries numerical relation.  De Vries remains a separate source-action
problem involving a named operator, a retained response variable, and
canonical normalization.

\section{The Three-Dimensional Prototype}

The usual custodial picture can be read from the Hopf fibration
\[
S^1\longrightarrow S^3\longrightarrow \CP^1 .
\]
For the Standard Model Higgs,
\[
H\in \CC^2\simeq \RR^4,\qquad |H|=\frac{v}{\sqrt 2},
\]
the fixed-radius vacuum orbit is \(S^3\).  The group acting on the ambient
real four-space is
\[
SO(4)\simeq \frac{SU(2)_L\times SU(2)_R}{\mathbb Z_2}.
\]
The diagonal subgroup after the vacuum choice is the custodial group.

The mass relation uses the ambient quadratic form on \(\RR^4\), not only the
topology of the orbit.  With a single doublet vacuum
\[
H_0=\begin{pmatrix}0\\ v/\sqrt2\end{pmatrix},
\]
the charged and neutral stiffnesses match:
\[
\kappa_c=\kappa_n,\qquad
\rho=\frac{m_W^2}{m_Z^2\cos^2\theta_W}=1 .
\]

\section{The Seven-Dimensional Carrier}

The CP2/Aloff--Wallach carrier is
\[
SO(3)\longrightarrow X_{1,1}=SU(3)/U(1)_{1,1}\longrightarrow \CP^2,
\]
where
\[
U(1)_{1,1}=\{\operatorname{diag}(z,z,z^{-2})\}.
\]
Equivalently, in the Wilking description,
\[
X_{1,1}\simeq (SU(3)\times SO(3))/U(2)_\Delta .
\]

At the base point the tangent representation splits as
\[
\mathfrak m=\mzero\oplus\mthree,
\]
with
\[
\mzero\simeq\so(3),\qquad \dim_{\RR}\mzero=3,
\]
and
\[
\mthree\simeq\CC^2\simeq\RR^4,\qquad \dim_{\RR}\mthree=4 .
\]
The invariant metric has two scales,
\[
g=x_1(-B)\big|_{\mzero}+x_2(-B)\big|_{\mthree}.
\]
The vertical scale \(x_1\) controls the \(SO(3)\) fibre.  The horizontal
scale \(x_2\) controls the Higgs four-block.

\section{Why \texorpdfstring{$\mathfrak m_3$}{m3} Is the Doublet}

The complexified \(\mthree\) representation has the charged split
\[
(\mthree)_{\CC}\simeq \mathbf 2_{+3}\oplus \mathbf 2_{-3}.
\]
With the electroweak normalization
\[
Y=\frac{q_8}{6},
\]
the summand \(\mathbf 2_{+3}\) has hypercharge \(Y=+1/2\).  Hence the
retained Higgs field is
\[
H\in \mathbf 2_{+3}\simeq \CC^2_{1/2}.
\]
The conjugate doublet belongs to \(\mathbf 2_{-3}\):
\[
\widetilde H=i\sigma_2 H^*\in \CC^2_{-1/2}.
\]

In a gauge-Higgs reading, use a higher electroweak carrier \(SU(3)_W\) and
embed the same \(U(1)_8\).  The centralizer is
\[
C_{SU(3)_W}(U(1)_8)=U(2)_W .
\]
The adjoint decomposes as
\[
\su(3)_W=(\mathbf 3_0\oplus\mathbf 1_0)
\oplus(\mathbf 2_{+3}\oplus\mathbf 2_{-3}).
\]
The internal connection components in the broken
\(SU(3)_W/U(2)_W\) directions match the charged \(\mthree\) block and give
one electroweak doublet.

\section{The Mass-Matrix Calculation}

Keep the invariant metric factor on \(\mthree\):
\[
g\big|_{\mthree}=x_2(-B)\big|_{\mthree}.
\]
The Higgs kinetic term has the form
\[
\mathcal L_H=
\kappa\, g_{\mthree}(D_\mu H,D^\mu H)-V(H,\Phi_0,\ldots).
\]
Use the usual electroweak covariant derivative
\[
D_\mu H
=\partial_\mu H
-i g W_\mu^a\frac{\sigma_a}{2}H
-i g' B_\mu\frac{1}{2}H .
\]
At
\[
H_0=\begin{pmatrix}0\\ v/\sqrt2\end{pmatrix},
\]
the charged terms give
\[
m_W^2=\frac{\kappa x_2 g^2v^2}{4}
\]
up to the common generator-normalization convention used for \((-B)\).

The neutral quadratic form is
\[
\frac{\kappa x_2v^2}{8}
\begin{pmatrix}W_\mu^3&B_\mu\end{pmatrix}
\begin{pmatrix}
g^2&-gg'\\
-gg'&g'^2
\end{pmatrix}
\begin{pmatrix}W^{3\mu}\\ B^\mu\end{pmatrix}.
\]
The zero eigenvector is the photon.  The massive eigenvalue gives
\[
m_Z^2=\frac{\kappa x_2(g^2+g'^2)v^2}{4}.
\]
Since
\[
\cos^2\theta_W=\frac{g^2}{g^2+g'^2},
\]
the common internal factor cancels:
\[
\rho=\frac{m_W^2}{m_Z^2\cos^2\theta_W}=1 .
\]

Equivalently, in stiffness notation,
\[
\kappa_c=\kappa_n .
\]
In the local normalization used elsewhere in the project, the common value
can be written as
\[
\kappa_c=\kappa_n=\frac{3\kappa v^2x_2}{4},
\]
with the same conclusion for \(\rho\).

\section{Schur's Lemma as the Geometric Reason}

The mass calculation above is the familiar one-doublet computation.  The
geometric reason it survives the Aloff--Wallach carrier is Schur's lemma.
The block \(\mthree\) is one irreducible real \(U(2)\)-module.  Hence every
\(U(2)\)-invariant symmetric bilinear form on \(\mthree\) is a scalar
multiple of the standard one:
\[
g\big|_{\mthree}=x_2(-B)\big|_{\mthree}.
\]
No independent charged and neutral stiffness can appear inside the same
single \(\mthree\) block.

The vertical block \(\mzero\) has its own scale \(x_1\).  That scale affects
the fibre geometry, the squash/radion action, and KK tower data.  It does
not split the weak stiffnesses as long as the electroweak vacuum stays in
\(\mthree\).

\section{The Vertical Triplet Hazard}

The vertical sector can contain an adjoint-like scalar \(\Phi_0\) associated
with \(\mzero\).  A minimal effective potential has the schematic form
\[
V(H,\Phi_0)=
m_3^2 H^\dagger H+\lambda_3(H^\dagger H)^2
\frac{m_0^2}{2}\Tr(\Phi_0^2)
\lambda_{03}H^\dagger H\,\Tr(\Phi_0^2)+\cdots .
\]
The custodial branch uses
\[
\langle H\rangle\ne0,\qquad \langle\Phi_0\rangle=0 .
\]
After the doublet condenses, stability of the vertical triplet requires
\[
m_{0,\mathrm{eff}}^2+\frac{\lambda_{03}v^2}{2}>0 .
\]
A vertical triplet vacuum gives a different vector-mass pattern and spoils
the single-doublet \(\rho=1\) argument.  The triplet therefore supplies a
real dynamical test, not a harmless decoration.

\section{D6 and D7 Roles}

The D7 endpoint is \(X_{1,1}\).  It also fibres over the D6 flag phase:
\[
S^1=T^2/U(1)_{1,1}\longrightarrow X_{1,1}\longrightarrow
Y_6=SU(3)/T^2 .
\]
The D6 phase has
\[
S^2\longrightarrow Y_6\longrightarrow \CP^2 .
\]
It carries the branch-cover and spin-\(c\) diagnostic data used elsewhere in
the programme, including the branched covering
\[
\CP^2\longrightarrow S^4
\]
with branch locus \(\RP^2\).

For the custodial lecture, the role of D6 is branch selection.  D6 defect,
boundary, or CSDR/Yang--Mills terms may shift
\[
m_{3,\mathrm{eff}}^2,\qquad m_{0,\mathrm{eff}}^2,\qquad \lambda_3,\qquad
\lambda_{03}.
\]
Those terms help only if they select the doublet branch and preserve the
isotropic real metric on \(\mthree\simeq\RR^4\).

The D7 gravitational action controls the radion and the squash ratio
\[
a=\frac{x_1}{x_2}.
\]
It can stabilize or destabilize the compactification shape.  The tree-level
\(\rho=1\) calculation depends on the single scalar \(x_2\) on \(\mthree\),
so \(a\) enters the custodial relation only through possible changes in the
effective potential that selects the vacuum.

\section{Comparison with \texorpdfstring{$\CP^2\times S^3$}{CP2 x S3}}

The product geometry
\[
\CP^2\times S^3
\]
is a useful comparison object.  The round \(S^3\) can model the fixed-radius
orbit of a doublet.  The product, however, has separate CP2 and \(S^3\)
factors, while \(X_{1,1}\) ties the \(SO(3)\) fibre and CP2 base through the
homogeneous Aloff--Wallach structure.  The split
\[
\mathfrak m=\mzero\oplus\mthree
\]
with \(\mthree\simeq\CC^2\) provides the CP2-indexed doublet block used in
the custodial argument.

\section{Brane Scan, Division Algebras, and the D6 Clue}

The relevant memory is the brane-scan and division-algebra pattern, not a
literal triality statement in both six and ten dimensions.  Super Yang--Mills
and the Green--Schwarz superstring single out
\[
D=3,4,6,10,
\]
corresponding to the normed division algebras
\[
\RR,\quad \CC,\quad \HH,\quad \mathbb O
\]
of dimensions
\[
1,\quad 2,\quad 4,\quad 8 .
\]
In Evans' formulation, division algebras explain the auxiliary-field
structures for super Yang--Mills in \(d=3,4,6,10\).\footnote{J. M. Evans,
\emph{Auxiliary Fields for Super Yang--Mills from Division Algebras},
arXiv:\href{https://arxiv.org/abs/hep-th/9410239}{hep-th/9410239}.}
Baez and Huerta give a modern self-contained account: nonabelian
Yang--Mills fields minimally coupled to massless spinors are supersymmetric
precisely in these spacetime dimensions, because the corresponding
trilinear spinor identity is controlled by the normed division
algebras.\footnote{J. C. Baez and J. Huerta,
\emph{Division Algebras and Supersymmetry I},
arXiv:\href{https://arxiv.org/abs/0909.0551}{0909.0551}.}

Triality belongs most sharply to the octonionic corner.  In ten-dimensional
massless physics the little group is \(\operatorname{Spin}(8)\), where the
vector and two chiral spinor representations
\[
\mathbf 8_v,\qquad \mathbf 8_s,\qquad \mathbf 8_c
\]
are permuted.  Six dimensions occupy the quaternionic corner: the chiral
spinors are quaternionic, and the Fierz identities have the same
division-algebra origin, without the full \(\operatorname{Spin}(8)\)
triality.

For the current programme, the useful consequence is a D6 consistency clue.
A six-dimensional parent action with supersymmetric Yang--Mills ancestry can
fix relative bosonic coefficients by the same algebraic identities that make
the D6 theory consistent.  Such a source action may select the \(\mthree\)
doublet branch, suppress the \(\mzero\) triplet, or rule out a proposed D6
defect term.  It still has to compute any De Vries candidate data from the
same reduced action:
\[
\mu_r,\qquad \gamma(u)^2,\qquad \alpha,\beta,\qquad
c=\frac{\beta}{\alpha^2},\qquad \sigma_r^{\mathrm{can}} .
\]
The division-algebra structure can enter as a source-action constraint.  It
cannot replace the source-action calculation.

\section{Separation from the De Vries Problem}

The custodial result proven here is
\[
\mthree\simeq\CC^2\simeq\RR^4,\qquad
\langle H\rangle\in\mthree,\qquad
\kappa_c=\kappa_n,\qquad
\rho=1 .
\]
The active De Vries goal asks for a different object: a source action that
computes a named operator, its spectral multiplier, the response kinetic
scale, and the signed Hessian coefficient in canonical variables.  The
custodial calculation supplies a stable electroweak control condition.  It
does not authorize inserting the desired spectral formula as an input.

\section{Checklist for Use in a Bigger KK Paper}

A larger paper can use this lecture as the custodial chapter once the
following items are stated explicitly:

\begin{enumerate}[leftmargin=2em]
  \item The geometry is \(SO(3)\to X_{1,1}\to\CP^2\) with
  \(X_{1,1}=SU(3)/U(1)_{1,1}\).
  \item The tangent representation splits as
  \(\mathfrak m=\mzero\oplus\mthree\), with \(\dim\mzero=3\) and
  \(\dim\mthree=4\).
  \item The electroweak doublet is the \(\mathbf 2_{+3}\) summand of
  \((\mthree)_{\CC}\), with \(Y=q_8/6\).
  \item The vacuum lies in the single \(\mthree\) doublet branch.
  \item The vertical triplet satisfies
  \(\langle\Phi_0\rangle=0\) and a positive effective mass condition.
  \item Endpoint, defect, and boundary terms preserve the isotropic
  quadratic form on \(\mthree\), or their correction to \(\rho\) is computed.
  \item De Vries material appears only after a source-action computation
  supplies the operator scale and Hessian coefficient.
\end{enumerate}

\section{Local Source Map}

The following project files contain the longer derivations and audits behind
this lecture:

\begin{itemize}[leftmargin=2em]
  \item \texttt{out/CUSTODIAL\_3D\_TO\_7D\_EXTENSION.md}: Hopf-to-Aloff--Wallach
  custodial dictionary.
  \item \texttt{out/KK\_HIGGSING\_CUSTODIAL\_MODEL.md}: gauge-Higgs model with
  \(SU(3)_W\to U(2)_W\) and the matched \(\mthree\) doublet.
  \item \texttt{out/KK\_HIGGS\_POTENTIALS\_1D3D\_5D7D.md}: D6 and D7 potential
  terms, branch conditions, and product comparison.
  \item \texttt{out/GRAVITY\_ACTION\_D765\_D321.md}: KK metric and
  Einstein--Hilbert decomposition across D7/D6/D5 and D3/D2/D1.
  \item \texttt{notes/Physics\_First\_KK\_Lagrangian\_Hessian\_Audit.md}:
  separation between ordinary KK Higgs physics and the De Vries
  source-action problem.
\end{itemize}

\end{document}
